% Part: first-order-logic
% Chapter: models-theories
% Section: expressing-relations

\documentclass[../../../include/open-logic-section]{subfiles}

\begin{document}

\olfileid{fol}{mat}{exr}

\olsection{التعبير عن العلاقات في \article{structure}
  \printtoken{S}{structure}}

\begin{explain}
من أعمال ال!!{formula}s أن تؤدي خواص وعلاقات في
!!a{structure}~$\Struct M$ بأوليات لغة~$\Lang L$
التي لها~$\Struct M$. وبيان ذلك أن !!{domain}~$\Struct M$
مجموعة من الكائنات، وأن مدلولات ال!!{constant}s وال!!{function}s
وال!!{predicate}s في~$\Struct M$، على الترتيب، كائنات من
$\Domain \OLId{latin-looped}{M}$، ودوال على~$\Domain \OLId{latin-looped}{M}$، وعلاقات على~$\Domain \OLId{latin-looped}{M}$.
فإذا كان $\Obj A^\OLMathNumeral{2}_\OLMathNumeral{0}$ من $\Lang L$، عيّنت له
$\Struct M$ العلاقة~$\OLId{latin-looped}{R} = \Assign{{\Obj
  A^\OLMathNumeral{2}_\OLMathNumeral{0}}}{\OLId{latin-looped}{M}}$. وحينئذ تعبر الصيغة $\Atom{\Obj A^2_0}{\Obj v_1, \Obj v_2}$
\emph{عن} تلك العلاقة بعينها. ومعنى ذلك أنه متى كان إسناد
المتغيرات~$\OLId{latin-ordinary}{s}$ يرسل $\Obj v_\OLMathNumeral{1}$ إلى $\OLId{latin-ordinary}{a} \in \Domain{\OLId{latin-looped}{M}}$ و$\Obj v_\OLMathNumeral{2}$
إلى $\OLId{latin-ordinary}{b} \in \Domain \OLId{latin-looped}{M}$، كان
\[
Rab \text{\quad iff\quad} \Sat{\OLId{latin-looped}{M}}{\Atom{\Obj A^2_0}{\Obj v_1, \Obj v_2}}[\OLId{latin-ordinary}{s}].
\]
ولا يستغنى هنا عن إسنادات المتغيرات؛ فلا يكفي قولنا
«$Rab$ إذا وفقط إذا كان $\Sat{\OLId{latin-looped}{M}}{\Atom{\Obj A^2_0}{a, b}}$»،
لأن $\OLId{latin-ordinary}{a}$ و$\OLId{latin-ordinary}{b}$ ليسا من رموز اللغة، بل !!{element}s
من~$\Domain{\OLId{latin-looped}{M}}$.

ولسنا محصورين في ال!!{formula}s الذرية؛ إذ تؤلّف الروابط
والمكممات منها !!{formula}s مركبة. وبها تُعرّف علاقات لا تكون
مكوّنات معيّنة مباشرة في~$\Struct M$. فالمطلوب بيان سبيل
هذا التعريف، وتمييز العلاقات التي يمكن تعريفها في !!a{structure}.
\end{explain}

\begin{defn}
لتكن $!A(\Obj v_\OLMathNumeral{1},\dots, \Obj v_\OLId{latin-ordinary}{n})$ !!a{formula} من $\Lang L$ لا
تقع فيها حرة إلا $\Obj v_\OLMathNumeral{1}$،\dots، $\Obj v_\OLId{latin-ordinary}{n}$، ولتكن $\Struct M$
!!a{structure} لـ~$\Lang L$. تعبر $!A(\Obj v_\OLMathNumeral{1},\dots, \Obj v_\OLId{latin-ordinary}{n})$
\emph{عن العلاقة}~$\OLId{latin-looped}{R} \subseteq \Domain \OLId{latin-looped}{M}^\OLId{latin-ordinary}{n}$ إذا وفقط إذا كان
\[
Ra_\OLMathNumeral{1}\dots \OLId{latin-ordinary}{a}_\OLId{latin-ordinary}{n} \text{\quad iff\quad} \Sat{\OLId{latin-looped}{M}}{\Atom{!A}{\Obj
    v_1,\dots, \Obj v_n}}[\OLId{latin-ordinary}{s}]
\]
لكل إسناد متغيرات~$\OLId{latin-ordinary}{s}$ يكون فيه $\OLId{latin-ordinary}{s}(\Obj v_\OLId{latin-ordinary}{i}) = \OLId{latin-ordinary}{a}_\OLId{latin-ordinary}{i}$ ($\OLId{latin-ordinary}{i} = \OLMathNumeral{1},
\dots, \OLId{latin-ordinary}{n}$).
\end{defn}

\begin{ex}
في النموذج القياسي للحساب~$\Struct N$، تعبر ال!!{formula} $\Obj
v_\OLMathNumeral{1} < \Obj v_\OLMathNumeral{2} \lor \eq[\Obj v_\OLMathNumeral{1}][\Obj v_\OLMathNumeral{2}]$ عن العلاقة $\le$ على
~$\Nat$. وتعبر ال!!{formula} $\eq[\Obj v_\OLMathNumeral{2}][\Obj v_\OLMathNumeral{1}']$ عن علاقة
الخلف، أي العلاقة $\OLId{latin-looped}{R} \subseteq \Nat^\OLMathNumeral{2}$ التي يصدق فيها $Rnm$ إذا كان
$\OLId{latin-ordinary}{m}$ خلف~$\OLId{latin-ordinary}{n}$. وتعبر الصيغة $\eq[\Obj v_\OLMathNumeral{1}][\Obj v_\OLMathNumeral{2}']$ عن علاقة السلف.
وتعبر ال!!{formula}s $\lexists[\Obj v_\OLMathNumeral{3}][(\eq/[\Obj v_\OLMathNumeral{3}][\Obj 0] \land
  \eq[\Obj v_\OLMathNumeral{2}][(\Obj v_\OLMathNumeral{1} + \Obj v_\OLMathNumeral{3})])]$ و$\lexists[\Obj
  v_\OLMathNumeral{3}][\eq[(\Obj v_\OLMathNumeral{1} + {\Obj v_\OLMathNumeral{3}}')][\Obj v_\OLMathNumeral{2}]]$ كلتاهما عن العلاقة
$\Obj <$. ومن ثم يُستغنى في أصل لغة الحساب عن رمز المحمول~$<$؛
إذ يردّ إلى الرموز الأخرى بالتعريف.
\end{ex}

\begin{explain}
ولا يختص هذا النظر ب!!{structure}s معيّنة؛ بل يعمّ كل استعمال
للغة في وصف نموذج مقصود أو نماذج مقصودة، أي في البحث عن النظريات.
فكثير من النظريات لا يجعل من ال!!{predicate}s الأولية إلا قليلًا،
مع أن مجال النظرية يشتمل على علاقات أخرى كثيرة لها شأن فيه.
فإذا أمكن أداء هذه العلاقات أداءً منتظمًا بالعلاقات التي تؤول
ال!!{predicate}s الأولية، قلنا إن تلك العلاقات تقبل
\emph{التعريف} في اللغة.
\end{explain}

\begin{prob}
ضع من !!{formula}s اللغة $\Lang L_\OLId{latin-looped}{A}$ ما يعرّف العلاقات الآتية:
\begin{enumerate}
\item $\OLId{latin-ordinary}{n}$ بين $\OLId{latin-ordinary}{i}$ و$\OLId{latin-ordinary}{j}$؛
\item $\OLId{latin-ordinary}{n}$ يقسم $\OLId{latin-ordinary}{m}$ دون باق (أي إن $\OLId{latin-ordinary}{m}$ مضاعف لـ~$\OLId{latin-ordinary}{n}$)؛
\item $\OLId{latin-ordinary}{n}$ عدد أولي (أي لا يقسم~$\OLId{latin-ordinary}{n}$ دون باق عدد غير $\OLMathNumeral{1}$ و$\OLId{latin-ordinary}{n}$).
\end{enumerate}
\end{prob}

\begin{prob}
لتكن الصيغة $!A(\Obj v_\OLMathNumeral{1}, \Obj v_\OLMathNumeral{2})$ تعبر عن العلاقة $\OLId{latin-looped}{R}
\subseteq \Domain \OLId{latin-looped}{M}^\OLMathNumeral{2}$ في !!a{structure}~$\Struct M$. أوجد صيغًا
تعبر عن العلاقات الآتية:
\begin{enumerate}
\item معكوس العلاقة $\OLId{latin-looped}{R}^{-\OLMathNumeral{1}}$؛
\item الضرب النسبي $\OLId{latin-looped}{R} \mid \OLId{latin-looped}{R}$؛
\end{enumerate}
أيمكن إيجاد عبارة تؤدي $\OLId{latin-looped}{R}^+$، وهو الإغلاق المتعدي لـ~$\OLId{latin-looped}{R}$؟
\end{prob}

\begin{prob}
لتكن $\Lang{L}$ لغةً رموزها غير المنطقية محصورة في محمول ثنائي المحل
$<$ (فلا !!{constant}s ولا !!{function}s ولا !!{predicate}s أخرى
فيها---باستثناء~$\eq$ بالطبع). ولتكن $\Struct{N}$ البنية التي فيها
$\Domain{\OLId{latin-looped}{N}} = \Nat$ و$\Assign{<}{\OLId{latin-looped}{N}} = \Setabs{\tuple{\OLId{latin-ordinary}{n},\OLId{latin-ordinary}{m}}}{\OLId{latin-ordinary}{n} <
  \OLId{latin-ordinary}{m}}$. أثبت ما يأتي:
\begin{enumerate}
\item $\{ \OLMathNumeral{0} \}$ قابلة للتعريف في $\Struct{N}$؛
\item $\{ \OLMathNumeral{1} \}$ قابلة للتعريف في $\Struct{N}$؛
\item $\{ \OLMathNumeral{2} \}$ قابلة للتعريف في $\Struct{N}$؛
\item لكل $\OLId{latin-ordinary}{n} \in \Nat$، المجموعة $\{ \OLId{latin-ordinary}{n} \}$ قابلة للتعريف في
  $\Struct{N}$؛
\item كل مجموعة جزئية منتهية من $\Domain{\OLId{latin-looped}{N}}$ قابلة للتعريف في
  $\Struct{N}$؛
\item كل مجموعة جزئية من $\Domain{\OLId{latin-looped}{N}}$ متممتها منتهية قابلة للتعريف
  في $\Struct{N}$ (حيث تكون متممة $\OLId{latin-looped}{X} \subseteq \Nat$ منتهية إذا
  وفقط إذا كانت $\Nat \setminus \OLId{latin-looped}{X}$ منتهية).
\end{enumerate}
\end{prob}


\end{document}
